% Auto-generated by python/build_tables.py
\begin{table*}[!htbp]
\centering
\begin{threeparttable}
\caption{Baseline parameterization of the simulation model.}
\label{tab:baseline-parameters}
\scriptsize
\begin{tabular}{p{0.13\textwidth}p{0.25\textwidth}p{0.12\textwidth}p{0.42\textwidth}}
\toprule
\textbf{Parameter} & \textbf{Meaning} & \textbf{Value} & \textbf{Managerial rationale} \\
\midrule
$n_U$ & University agents & 20.0 & Symmetric modular R and D collaboration setting. \\
$n_I$ & Industry agents & 20.0 & Symmetric modular R and D collaboration setting. \\
$p_{in}$ & Internal tie probability & 0.2 & Within-module interaction is more frequent than cross-boundary interaction. \\
$p_{out}$ & Baseline cross-boundary tie probability & 0.03 & Cross-boundary ties are sparse in unmanaged modular collaboration. \\
$k$ & Cross-boundary tie budget & 12.0 & Fixed scarce interface capacity in controlled architectures. \\
$b$ & Boundary spanners per side & 2.0 & Concentrated boundary-spanning role allocation. \\
$w_0$ & Initial relational confidence & 0.4 & Initial relations are neither absent nor ready. \\
$\alpha$ & Reinforcement rate & 0.08 & Successful interactions increase relational confidence. \\
$\beta$ & Decay rate & 0.02 & Failed interactions erode relational confidence. \\
$\pi_{in}$ & Internal success probability & 0.8 & Interactions within the same module are easier. \\
$\pi_{out}$ & Ordinary cross-boundary success probability & 0.55 & Cross-boundary interactions are more difficult. \\
$\pi_{BS}$ & Boundary-spanning success probability & 0.65 & Boundary spanners have a modest translation advantage. \\
$\theta$ & Tie-level readiness threshold & 0.8 & A tie must reach high relational confidence to count as ready. \\
$q$ & Boundary-level readiness threshold & 0.8 & The boundary is ready when most cross-boundary ties are ready. \\
$T_{\max}$ & Simulation horizon & 50000.0 & Maximum number of interaction events per trajectory. \\
$NG$ & Graph realizations & 50.0 & Independent network structures per condition. \\
$NT$ & Trajectories per graph & 50.0 & Stochastic interaction histories per network structure. \\
\bottomrule
\end{tabular}
\begin{tablenotes}
\footnotesize
\item Notes: Parameter values define a transparent reference case rather than an empirical calibration. Robustness analyses vary the main parameters linked to translation capability, boundary-spanner load, interaction selection, and readiness thresholds.
\end{tablenotes}
\end{threeparttable}
\end{table*}
